% ACM conference typography for a course assignment.
\documentclass[sigconf,nonacm]{acmart}

% Wrap database queries within the ACM column without changing their syntax.
\usepackage{listings}
\lstdefinestyle{searchquery}{
  basicstyle=\small\ttfamily,
  columns=fullflexible,
  keepspaces=true,
  breaklines=true,
  breakatwhitespace=false,
  showstringspaces=false,
  upquote=true
}

\begin{document}

\title{Q-Matrix}

\author{BG Rodrigo L Abellanosa}
\email{babellanosa3@gatech.edu}

\author{Christopher J Brookhart}
\email{cbrookhart3@gatech.edu}

\author{Friedrich K Eibl}
\email{feibl3@gatech.edu}

\author{John A McCloskey}
\email{jmccloskey30@gatech.edu}

\author{Monroe Stephenson}
\email{mstephenson31@gatech.edu}

\affiliation{%
  \institution{Georgia Institute of Technology}
  \city{Atlanta}
  \state{Georgia}
  \country{USA}
}

\renewcommand{\shortauthors}{Abellanosa et al.}

\begin{abstract}
A Q-matrix is a binary matrix linking assessment items to the latent skills or attributes they are intended to measure. Cognitive diagnostic models use these relationships to interpret responses, so errors in the matrix can produce misleading assessments. Q-matrices have traditionally been constructed through expert judgment, while statistical and machine learning methods provide ways to estimate or revise these relationships. Comparing these methods across domains can help determine which are suitable for particular data and assessment requirements.

This systematic literature review will examine research from the past ten years on Q-matrix generation, validation, and refinement in domains such as educational measurement, intelligent tutoring systems, and psychological assessment. We will compare methodological developments and examine their reported strengths, limitations, and the problems they address. We will also investigate how methods perform under different data conditions and how that performance is measured. Structured searches and consistent screening and data extraction criteria will guide the review. Comparisons will consider sample size, attribute structure, and interpretability requirements. The goal is to assess which methodological choices are supported by the published evidence and where differences in data or evaluation practices limit comparisons across fields.

\end{abstract}

\keywords{Q-Matrix}

\maketitle

\section{Research Questions}

\subsection{Research Question 1}

Which methods for generating, validating, and refining Q-matrices have been proposed or applied in peer-reviewed studies from 2021 to 2026, and how does the percentage of studies using each method family differ across educational measurement, intelligent tutoring systems and learning analytics, and psychological and clinical assessment?

Overview. Q-Matrix methods have developed in different domains, with little cross-domain research conducted to date. RQ1 will facilitate a more holistic understanding of the Q-Matrix methods.

\paragraph{Novelty and Differentiation.} Based on an initial research assessment (manual searches of Scopus, Web of Science Core Collection, APA PsycInfo, and ACM Digital Library with an AI review of results, along with an AI-conducted search of openly accessible sources), it appears this question is novel. Existing identified reviews (Qin et al., 2026; Wang et al., 2026; Coates, 2026) either focus exclusively on cognitive diagnostic models or focus on a single domain.

\paragraph{Relevance and Intellectual Merit/Broader Impact.} Different domains rarely cite each other. Further, LLM-based work has grown approximately 4-fold from 2024-2026. Each of these facts suggests an opportunity to facilitate cross-domain understanding in a changing field.

\paragraph{Specificity and Measurability.} RQ1 focuses on a specific population of peer-reviewed studies from 2021-2026. Each study will be coded by:

\begin{itemize}
    \item task (e.g., generation, validation, or refinement)
    \item method family (e.g., expert-driven, empirical estimation such as learning factor analysis, regularization-based, machine-learning, or LLM-based approaches)
    \item domain (e.g., educational measurement; intelligent tutoring systems and learning analytics; psychological and clinical assessment)
    \item publication year
\end{itemize}
Output will include percentage-based conclusions.

\paragraph{Answering the Research Question through Literature Search.} The relevant population of papers from applicable databases (e.g., Scopus, Web of Science Core Collection, APA PsycInfo, and ACM Digital Library – see Part 2 question 1 for more detail) will be extracted and then categorized by method family, domain, and task (e.g., generate, validate, or refine). Results will include a family-by-domain table with counts and percentages.

\paragraph{Feasibility.} An AI-facilitated initial scoping search of openly accessible sources and review of sources from manual database searches found 963 records (after de-duplication), in line with the scope suggested in the Group Project Check In 1 assignment description.

\subsection{Research Question 2}

What data conditions (sample size, attribute structure, interpretability standards) and evaluation criteria (metrics and benchmarks) are reported for each method family in each domain, and which evaluation criteria are most common across domains?

Overview. RQ2 seeks to answer what evidence studies of Q-Matrices rely on and how success is measured for them both within domains and across them.

\paragraph{Novelty and Differentiation.} Based on our initial research assessment, it appears the novelty of this question rests in its cross-domain focus. No review was found that compares data conditions and evaluation criteria across domains or asks which criteria domains share the most.

\paragraph{Relevance and Intellectual Merit/Broader Impact.} This review will identify where evidence is thinnest and those evaluation criteria that practitioners may find most generally applicable. It will enable practitioners to assess whether a method was applied to data like theirs before trusting the diagnosis the method produces in their testing.

\paragraph{Specificity and Measurability.} Every study will be coded to its data conditions (such as sample size, attribute structure, and interpretability standards). Its evaluation criteria will also be coded (e.g., the metrics and benchmarks applied to the Q-matrix). This will enable (i) the production of a table of data conditions by method family and domain and ii) a catalog of evaluation criteria identified and the domains in which they are used.

\paragraph{Answering the Research Question through Literature Search.} RQ2 will use the same studies as RQ1. The data conditions and evaluation criteria from the studies will be extracted directly from the papers.

\paragraph{Feasibility.} The scope identified by the initial search (963 papers) is feasible for human review for a group of 5 across a semester project. Automated extraction techniques can further facilitate feasibility.

\subsection{Research Question 3}

What strengths and limitations do authors report for each method family, what problem is each family designed to solve, and which limitations recur across studies and domains?

\paragraph{Novelty and Differentiation.} As with RQ1 and RQ2, existing reviews lack a cross-domain focus. Additionally, no identified existing review compares the strengths and weaknesses of the complete set of Q-Matrix methods contemplated by this review.

\paragraph{Relevance and Intellectual Merit/Broader Impact.} RQ3 separates limitations specific to individual domains and those that apply more broadly, which will help demonstrate where new methods are needed.

\paragraph{Specificity and Measurability.} Each study will be coded to the stated problem the method is intended to address and any strengths and limitations stated in the study.

\paragraph{Answering the Research Question through Literature Search.} RQ3 will focus on the same set of studies as RQ1 and RQ2. Strengths and limitations will be extracted from the introduction and discussion sections of each paper. Following extraction, thematic analysis will categorize the identified strengths and weaknesses, which will then be counted by family and domain.

\paragraph{Feasibility.} RQ3 is the most challenging of the three. Automated extraction methods are less feasible than the other RQ’s. However, thematic coding is a common qualitative method, so the main feasibility challenge is the limited time to complete the review. If it becomes apparent that extraction and thematic coding for RQ3 is infeasible within the semester, we will limit RQ3 coding to a stratified sample of studies by method family.

\section{Procedure}

\subsection{Sets of Databases and Search Strategy}

To capture the interdisciplinary nature of Q-matrix research, we will be looking at papers from the following databases: Scopus, Web of Science Core Collection, APA PsycInfo, ERIC, PubMed, ACM Digital Library, and IEEE Xplore. While Scopus and Web of Science provide literature from many different disciplines, PubMed and PsycInfo focus on the medical fields, ERIC on educational research, and ACM as well as IEEE on its applications in technology and computer science. This combination supports the cross-domain focus of our research questions. Searches will primarily focus on peer-reviewed studies published over the span of 2021-2026.

\subsubsection{Scopus (Advanced search)}

\begin{lstlisting}[style=searchquery]
( TITLE-ABS-KEY ( ( "Q-matri*" OR "Q matri*" ) AND ( "cognitive diagnos*" OR "diagnostic classification" OR "diagnostic assessment" OR DINA OR attribute* OR skill* OR "knowledge component*" OR assessment* ) ) OR TITLE-ABS-KEY ( ( "knowledge component*" OR "skill model*" OR "cognitive model discover*" OR "knowledge concept tagging" OR "skill tagging" OR "knowledge tagging" ) AND ( discover* OR refin* OR generat* OR tagging OR "data-driven" OR "learning factors analysis" OR automat* OR extract* ) ) ) AND PUBYEAR > 2020 AND PUBYEAR < 2027
\end{lstlisting}

\subsubsection{Web of Science Core Collection (Advanced search, pasted into the Query Preview box)}

\begin{lstlisting}[style=searchquery]
TS=((("Q matri*") AND ("cognitive diagnos*" OR "diagnostic classification" OR "diagnostic assessment" OR DINA OR attribute* OR skill* OR "knowledge component*" OR assessment*)) OR (("knowledge component*" OR "skill model*" OR "cognitive model discover*" OR "knowledge concept tagging" OR "skill tagging" OR "knowledge tagging") AND (discover* OR refin* OR generat* OR tagging OR "data driven" OR "learning factors analysis" OR automat* OR extract*)))
\end{lstlisting}

\subsubsection{APA PsycInfo (EBSCO advanced search).}

\begin{lstlisting}[style=searchquery]
( TI ( "cognitive diagnos*" OR "diagnostic classification model*" ) OR AB ( "cognitive diagnos*" OR "diagnostic classification model*" ) ) AND ( TI ( clinical OR psychiatr* OR psychopatholog* OR "mental health" OR symptom* OR DSM OR personality ) OR AB ( clinical OR psychiatr* OR psychopatholog* OR "mental health" OR symptom* OR DSM OR personality ) )
\end{lstlisting}

\subsubsection{ACM Digital Library (Advanced search, with the field set to Abstract).}

\begin{lstlisting}[style=searchquery]
"Q-matrix" OR "Q matrix"
\end{lstlisting}

and a second search:

\begin{lstlisting}[style=searchquery]
("knowledge component" OR "knowledge components" OR "skill tagging" OR "knowledge concept tagging") AND (generation OR discovery OR refinement OR tagging OR automated)
\end{lstlisting}

\subsubsection{IEEE Xplore (Advanced search-Command Search)}

\begin{lstlisting}[style=searchquery]
(("Q-matrix" OR "Q matrix") AND ("cognitive diagnosis" OR "cognitive diagnostic" OR "diagnostic classification" OR DINA)) OR (("knowledge component" OR "knowledge components" OR "skill tagging" OR "knowledge concept tagging") AND (generation OR discovery OR refinement OR tagging OR automated))
\end{lstlisting}

\subsection{Explanation of Search Strings}

\begin{lstlisting}[style=searchquery]
(  ("Q-matri*" OR "Q matri*")  AND  ("cognitive diagnos*" OR "diagnostic classification*"   OR "diagnostic assessment*" OR DINA   OR attribute* OR skill* OR "knowledge component*" OR assessment*))OR(  ("knowledge component*" OR "skill model*"   OR "cognitive model discover*"   OR "knowledge concept tagging"   OR "skill tagging" OR "knowledge tagging")  AND  (discover* OR refin* OR generat*   OR tagging OR "data driven"   OR "learning factors analysis"   OR automat* OR extract*))OR(  ("cognitive diagnos*" OR "diagnostic classification*")  AND  (clinical OR psychiatr* OR psychopatholog*   OR "mental health" OR symptom* OR DSM OR personality))
\end{lstlisting}

The search string will follow three complementary search paths. The first targets studies explicitly using Q-matrix terminology together with terms associated with cognitive diagnosis and assessment. The second captures related research that may describe the generation or refinement of skill and knowledge-component mappings without explicitly referring to them as Q-matrices. The third targets cognitive diagnostic research in psychological and clinical settings to support the cross-domain scope of the review. Within each group, OR combines synonymous or related terms, while AND requires records to contain concepts from both groups. Quotation marks are used for multi-word phrases, and the wildcard * captures multiple endings of a term, such as refin* for refine, refinement, and refining. The exact syntax will be adapted to each database's search functionality while preserving these conceptual groups. Some terms will be more relevant for some databases than others, and we might also expand on these terms.

\subsection{Inclusion and Exclusion Criteria with Explanation}

\textbf{Exclusion Criteria.} Studies will be included if they were published in peer-reviewed English conferences or journals between 2021 and 2026 and address the generation, estimation, validation, correction, or refinement of a Q-matrix for mapping of certain skills or attributes. In addition, the studies need to provide sufficient information to contribute to at least one of our predefined research questions.

\textbf{Inclusion Criteria.} Studies will be excluded, if they use the term Q-matrix in an unrelated context, were published before our predefined review period, have not yet received successful peer review, or are not available in English. We will also aim to exclude any duplicates of the same publications or any studies that do not contribute to our research question in a meaningful way. Lastly, we will not include any publications that exclusively take advantage of pre-existing Q-Matrices and do not contribute to their refinement or validation in any way, as our review is specifically about the generation of Q-Matrices.

We will aim to deploy these criteria both for initial abstract-level screening as well as further paper-level screening.

\subsection{What Information Will Be Extracted to Answer the RQs}

For each paper included in the final review, we will extract the information needed to answer the three research questions. Basic information such as publication year and application domain will be recorded first. Each study will then be classified according to the Q-matrix task it addresses, such as generation, validation, or refinement.

To answer RQ1, we will record the method or methods used in each study and classify them into broader methodological categories such as expert-driven approaches, empirical or statistical estimation methods, regularization-based methods, machine-learning or deep-learning approaches, and LLM-based approaches. We will also record the application domain of each study, including educational measurement, intelligent tutoring systems and learning analytics, and psychological or clinical assessment. These data will allow us to compare the number and percentage of studies using each method family across domains and publication years.

To answer RQ2, we will extract information about the data conditions under which each method was studied or applied. This will include sample size, number of items, number of attributes or skills, type of data used, whether the study used simulated or empirical data, and characteristics of the attribute structure. We will also record any interpretability requirements or considerations explicitly discussed by the authors. In addition, we will extract the evaluation criteria used to assess Q-matrix performance, including metrics, benchmarks, comparison methods, and validation procedures. These may include measures such as Q-matrix recovery, classification accuracy, agreement with expert judgments, model-fit statistics, predictive performance, or other criteria reported by the authors.

To answer RQ3, we will extract the problem that each method is intended to address, together with the strengths and limitations of the method as reported by the authors. These statements will be grouped into recurring themes so that common strengths, limitations, and unresolved problems can be compared across methodological families and application domains.

If a study does not report a particular item of information, that variable will be recorded as not reported rather than inferred by the reviewers.

\subsection{Current Estimate for Amount of Work}

Using the current search strings, the 2021-2026 publication range, and the databases included in the review, our initial searches produced 963 records. These 963 records will be subjected to screening and filtering rather than the final number of records expected to receive full review and data extraction. The majority will be screened using titles and abstracts according to the inclusion and exclusion criteria. Records that are potentially relevant will be subjected to full-text screening and only studies that satisfy the final eligibility criteria will be included in the data-extraction stage.

With five members in the group, this corresponds to approximately 193 title and abstract reviews per member if divided evenly. This number should get smaller after initial screening.

\begin{thebibliography}{3}

\bibitem{qin2026}
Qin, Haijiang, et al. “From Theory to Practice: A Comprehensive Toolkit for Q-Matrix Validation in Cognitive Diagnosis.” Behavior Research Methods, vol. 58, no. 6, May 2026, p. 151. Springer Link, \href{https://doi.org/10.3758/s13428-026-02991-5}{https://doi.org/10.3758/s13428-026-02991-5}.

\bibitem{coates2026}
Coates, Adam. “A Systematic Review of the Implications of Diagnostic Classification Modeling for Cognitive Models of Mathematics.” Review of Educational Research, Sept. 2026, p. 00346543261480692. DOI.org (Crossref), \href{https://doi.org/10.3102/00346543261480692}{https://doi.org/10.3102/00346543261480692}.

\bibitem{wang2026}
Wang, Chun, et al. “Review of Cognitive Diagnostic Models ( CDMs ): Recent Methodological Advancements for Addressing Practical Challenges.” British Journal of Mathematical and Statistical Psychology, July 2026, p. bmsp.70066. DOI.org (Crossref), \href{https://doi.org/10.1111/bmsp.70066}{https://doi.org/10.1111/bmsp.70066}.

\end{thebibliography}

\end{document}
